\section{Related Work}
\label{sec:related}

\paragraph{Recognising gestures from pixels.} Early systems found the hand
by its skin colour or by removing the background, and then described its shape;
Rautaray and Agrawal give a survey~\cite{rautaray2015survey}. These methods
break easily when the light or the background changes. Deep learning changed
the field. 3D CNNs with recurrent layers recognise moving gestures in
video~\cite{molchanov2016online}, and two-stage CNN systems detect and
classify gestures in real time~\cite{kopuklu2019realtime}. Large datasets
such as Jester~\cite{materzynska2019jester} and HaGRID~\cite{kapitanov2024hagrid}
made this possible. These models work very well, but they read raw pixels, so
they need a lot of training data and usually a GPU.

\paragraph{Finding hand points.} Another line of work finds the joints of
the hand instead of classifying pixels. Simon \etal~\cite{simon2017hand} detect
21 hand keypoints in a single image, and MediaPipe
Hands~\cite{zhang2020mediapipe} does this in real time on a phone CPU. Once we
have these 21 points, the background, skin colour and lighting are already
gone, and the gesture problem becomes much smaller: only 63 numbers per frame.
This is why we build on MediaPipe.

\paragraph{Virtual mouse programs.} Many open-source projects join MediaPipe
with PyAutoGUI~\cite{pyautogui}. They count fingers as ``up'' or ``down'' with
fixed thresholds and move the cursor with the fingertip. They are easy to
build, but, as said above, they are fragile, shaky and click by mistake.
They are also rarely measured: they show a demo, not numbers.

\paragraph{Smoothing and measuring a pointer.} Shaky input is usually
smoothed with a moving average or a Kalman filter~\cite{kalman1960}. Both
trade shake for delay. The One-Euro filter~\cite{casiez2012oneeuro} changes
its smoothing with the speed of the hand, which removes most of this
trade-off. Pointing devices are usually compared with Fitts'
law~\cite{fitts1954,soukoreff2004}.

\paragraph{Where our work fits.} We take the keypoint idea from the second
line of work and the virtual-mouse goal from the third, and we add what is
missing in between: a learned and position-free classifier tested on new
users, and a cursor that is designed and measured to be steady and reliable.
